% =============================================================
%  chapter5_photo_album.tex -- AUTO-GENERATED, DO NOT EDIT BY HAND.
%  Appendix A photo album: the photos in ch5_photo_refs.APPENDIX_PHOTOS,
%  two large photos per page, in numerical order, landscape photos first
%  and portrait photos after them. Written by
%  Chapter5/Figures/write_chapter5_photo_album.py from
%  Workspace/Draaikraal Photo Captions.docx; rerun it after editing the
%  captions or the selection. The other photos are in
%  chapter5_photo_supplement_album.tex (Chapter5/chapter5_photo_supplement.tex).
% =============================================================
\providecommand{\phantomsection}{}
\providecommand{\ChFivePhoto}[1]{#1}
% Two photos per page, photos at full size: landscape photos at the full text
% width (at most \ChFiveBigMax tall), portrait photos at 0.48 of it (at most
% \ChFiveTallMax tall). The captions (single-spaced) take the room that is left:
% TeX tries \small, \footnotesize, \scriptsize in turn and keeps the largest size
% at which both captions fit (the same size for both captions of a page). Only
% if even \scriptsize does not fit are the photos shrunk, just enough to make room.
% Each trial sets the whole page in a box and measures it.
\newsavebox\ChFiveImgA
\newsavebox\ChFivePageBox
\newsavebox\ChFiveImgB
\newsavebox\ChFiveCapA
\newsavebox\ChFiveCapB
\newlength\ChFiveImgH
\newlength\ChFiveImgW
\newlength\ChFiveCapW
\newif\ifChFiveFits
\newcommand{\ChFiveBigMax}{0.45\textheight}
\newcommand{\ChFiveTallMax}{0.62\textheight}
\newcommand{\ChFiveCapFont}{\small}
\newcommand{\ChFiveTotal}[1]{\dimexpr\ht#1+\dp#1\relax}
\newcommand{\ChFiveSetImg}[2]{%  {box}{image}
  \sbox#1{\includegraphics[width=\ChFiveImgW,height=\ChFiveImgH,keepaspectratio]{#2}}}
\newcommand{\ChFiveSetCaps}{%
  \sbox\ChFiveCapA{\parbox[t]{\ChFiveCapW}{\def\baselinestretch{1}\ChFiveCapFont\raggedright
    \setlength{\parskip}{2pt}\ChFiveCapTextA}}%
  \sbox\ChFiveCapB{\parbox[t]{\ChFiveCapW}{\def\baselinestretch{1}\ChFiveCapFont\raggedright
    \setlength{\parskip}{2pt}\ChFiveCapTextB}}}
\newcommand{\ChFiveCapsHeight}{% stacked (landscape) or side by side (portrait)
  \ifChFiveSideBySide
    \ifdim\ChFiveTotal\ChFiveCapA>\ChFiveTotal\ChFiveCapB \ChFiveTotal\ChFiveCapA
    \else\ChFiveTotal\ChFiveCapB\fi
  \else\dimexpr\ChFiveTotal\ChFiveCapA+\ChFiveTotal\ChFiveCapB\relax\fi}
\newif\ifChFiveSideBySide
\newcommand{\ChFiveTryFont}[1]{%  fits = the finished page (\ChFivePage) is no taller than the text block
  \ifChFiveFits\else
    \renewcommand{\ChFiveCapFont}{#1}\ChFiveSetCaps
    \setbox\ChFivePageBox\vbox{\ChFivePage}%
    \ifdim\dimexpr\ht\ChFivePageBox+\dp\ChFivePageBox+\topskip+2pt\relax<\textheight
      \ChFiveFitstrue\fi
  \fi}
\newcommand{\ChFiveFit}{%
  \ChFiveFitsfalse
  \ChFiveTryFont{\small}\ChFiveTryFont{\footnotesize}%
  \ChFiveTryFont{\scriptsize}}
\newcommand{\ChFiveBigOne}[4]{%  {width}{n}{image box}{caption box}
  \begin{minipage}[t]{#1}\centering
  \phantomsection\label{ch5-photo-#2}%
  \usebox#3\par
  \vspace{4pt}\raggedright\usebox#4\par
  \end{minipage}}
% Final check: the page is set in a trial box and, if it is still taller than
% the text block, both photos are shrunk by exactly the excess (a few passes).
\newcount\ChFiveIter
\newif\ifChFiveTooTall
\newlength\ChFiveExcess
\newcommand{\ChFiveSettle}[3]{%  {page content}{img1}{img2}
  \ChFiveIter=0
  \loop
    \setbox\ChFivePageBox\vbox{#1}%
    \setlength\ChFiveExcess{\dimexpr\ht\ChFivePageBox+\dp\ChFivePageBox-\textheight+\topskip+2pt\relax}%
    \ChFiveTooTallfalse
    \ifnum\ChFiveIter<6 \ifdim\ChFiveExcess>0pt \ChFiveTooTalltrue\fi\fi
  \ifChFiveTooTall
    \ifdim\ht\ChFiveImgA>\ht\ChFiveImgB
      \setlength\ChFiveImgH{\ht\ChFiveImgA}\else\setlength\ChFiveImgH{\ht\ChFiveImgB}\fi
    \ifChFiveSideBySide\else\setlength\ChFiveExcess{0.5\ChFiveExcess}\fi
    \setlength\ChFiveImgH{\dimexpr\ChFiveImgH-\ChFiveExcess-1pt\relax}%
    \ChFiveSetImg\ChFiveImgA{#2}\ChFiveSetImg\ChFiveImgB{#3}%
    \advance\ChFiveIter by 1
  \repeat}
\newcommand{\ChFiveReport}[2]{\typeout{ChFive album: photos #1 and #2, captions
  \expandafter\string\ChFiveCapFont\ifnum\ChFiveIter>0 \space (photos shrunk)\fi}}
% landscape-orientation photos: one above the other, full text width
\newcommand{\ChFiveWidePair}[6]{%  {n1}{img1}{cap1}{n2}{img2}{cap2}
  \ChFiveSideBySidefalse
  \setlength\ChFiveImgW{\linewidth}\setlength\ChFiveImgH{\ChFiveBigMax}%
  \setlength\ChFiveCapW{\linewidth}%
  \ChFiveSetImg\ChFiveImgA{#2}\ChFiveSetImg\ChFiveImgB{#5}%
  \def\ChFiveCapTextA{#3}\def\ChFiveCapTextB{#6}%
  \def\ChFivePage{\noindent\ChFiveBigOne{\linewidth}{#1}\ChFiveImgA\ChFiveCapA\par\vfill
    \noindent\ChFiveBigOne{\linewidth}{#4}\ChFiveImgB\ChFiveCapB\par}%
  \ChFiveFit
  \ChFiveSettle{\ChFivePage}{#2}{#5}\ChFiveReport{#1}{#4}%
  \ChFivePage\vfill\clearpage}
\newcommand{\ChFiveWideOne}[3]{%
  \setlength\ChFiveImgW{\linewidth}\setlength\ChFiveImgH{\ChFiveBigMax}%
  \setlength\ChFiveCapW{\linewidth}\renewcommand{\ChFiveCapFont}{\small}%
  \ChFiveSetImg\ChFiveImgA{#2}\def\ChFiveCapTextA{#3}\def\ChFiveCapTextB{}\ChFiveSetCaps
  \noindent\ChFiveBigOne{\linewidth}{#1}\ChFiveImgA\ChFiveCapA\par\clearpage}
% portrait-orientation photos: side by side, each 0.48 of the text width
\newcommand{\ChFiveTallPair}[6]{%
  \ChFiveSideBySidetrue
  \setlength\ChFiveImgW{0.48\linewidth}\setlength\ChFiveImgH{\ChFiveTallMax}%
  \setlength\ChFiveCapW{0.48\linewidth}%
  \ChFiveSetImg\ChFiveImgA{#2}\ChFiveSetImg\ChFiveImgB{#5}%
  \def\ChFiveCapTextA{#3}\def\ChFiveCapTextB{#6}%
  \def\ChFivePage{\noindent\ChFiveBigOne{0.48\linewidth}{#1}\ChFiveImgA\ChFiveCapA\hfill
    \ChFiveBigOne{0.48\linewidth}{#4}\ChFiveImgB\ChFiveCapB\par}%
  \ChFiveFit
  \ChFiveSettle{\ChFivePage}{#2}{#5}\ChFiveReport{#1}{#4}%
  \ChFivePage\clearpage}
\newcommand{\ChFiveTallOne}[3]{%
  \setlength\ChFiveImgW{0.48\linewidth}\setlength\ChFiveImgH{\ChFiveTallMax}%
  \setlength\ChFiveCapW{0.48\linewidth}\renewcommand{\ChFiveCapFont}{\small}%
  \ChFiveSetImg\ChFiveImgA{#2}\def\ChFiveCapTextA{#3}\def\ChFiveCapTextB{}\ChFiveSetCaps
  \noindent\ChFiveBigOne{0.48\linewidth}{#1}\ChFiveImgA\ChFiveCapA\par\clearpage}

\clearpage
\ChFiveWidePair{2}{Chapter5/Figures/generated/field_photos/DrBV2.jpg}{\textbf{Photo 2.} Decimeter scale lenticular inclusions of coarse-grained Anorthositic-Gabbro in magnetic, fine-grained Marginal Sill Phase. Clinopyroxene in inclusions is oikocrystic and up to 3 cm. Found alongside spheroidal pyroxenite aggregates of photo~\ChFivePhoto{4}.}{4}{Chapter5/Figures/generated/field_photos/DrBV4.jpg}{\textbf{Photo 4.} Pyroxenite Aggregates (colloquially referred to in the literature as “Tennis Balls”) 1-5 cm in diameter within fine-grained marginal zone gabbronorite. Pyroxenite aggregates have not been observed before this study outside of coarsely crystalline plutonic rocks, and only recently have authors argued that they could be hosted within the Marginal Sill Phase \citep{mitchell2023}. This study offers multiple unambiguous examples of pyroxenite aggregates hosted within fine-grained, magnetic, marginal sill phase.}
\ChFiveWidePair{49}{Chapter5/Figures/generated/field_photos/DrBV49.jpg}{\textbf{Photo 49.} A roadcut exposure of $\sim$30 cm’s of the UMA, the Upper Mottled Anorthosite, near the village at Draaikraal, south of the Draaikraal Lineament. The mottled anorthosite has an irregular, non-planar lower contact with a coarse-grained feldspathic pyroxenite. The UMA package has a characteristic assemblage of 3 distinct rocks types: Mottled anorthosite, a coarse-grained feldspathic pyroxenite, and a fine-grained green gabbronorite with blebby $<$2 cm clinopyroxene oikocrysts that weakly define modal layering (not pictured here). These distinct lithologies facilitate correlation with UMA to the north of the Draaikraal Lineament.}{59}{Chapter5/Figures/generated/field_photos/DrBV59.jpg}{\textbf{Photo 59.} A closeup of a laterite “ironstone” of unit IS. The ironstones are composed of 1-2 cm diameter goethite nodules often surrounding voids or porous-powdery limonite. The rock is cemented with a limonite-hematite matrix, and contains up to 20\% clasts of Steenkampsberg quartzite alluvium up to $\sim$ 1 cm in diameter. This rock type commonly hosted nests of ants.}
\ChFiveWidePair{61}{Chapter5/Figures/generated/field_photos/DrBV61.jpg}{\textbf{Photo 61.} A specimen containing interesting features commonly mapped as “Disturbed Gabbronorite” here, which generally refers to coarsely crystalline gabbronorites immediately stratigraphically above the deformed floor rocks of the Kliprivier Syncline. Large, Euhedral OPX grains are cemented with plagioclase and clinopyroxene oikocrysts, locally differentiated into bands and regions of pyroxenite or OPX phenocrysts in an anorthositic matrix, somewhat resembling the “porphyritic gabbronorite”. Such rocks may record the mechanical separation of pyroxenite and plagioclase to form pyroxenite spheroidal aggregates or “tennis balls”. The lens cap is 5.5 centimeters across.}{62}{Chapter5/Figures/generated/field_photos/DrBV62.jpg}{\textbf{Photo 62.} A freshly broken face of coarsely crystalline disturbed gabbronorite. Crystal size varies greatly throughout the hand sample, with some regions dominated by cm-scale crystals and some dominated by mm-scale crystals, suggesting perhaps that these rocks sustained channelized fluid flow in the finely crystalline regions after bulk crystallization. Even in a freshly broken face, the OPX faces appear to have a rusty discoloration. Several elongate feldspar inclusions are visible in the OPX phenocrysts.}
\ChFiveWidePair{65}{Chapter5/Figures/generated/field_photos/DrBV65.jpg}{\textbf{Photo 65.} Anastomosing cm-scale chromitite bands hosted within mottled anorthosite. The UG-1 layer consistently contains strings of chromitite hosted within mottled anorthosite just below the massive chromitite \citep{voordouw2009,mukherjee2017}. The most famous locality of UG-1 is at Dwars River. Here the oikocrystic pyroxenes are green when fresh, weathering to a rusty red color. A few isolated bands of chromite near the top are quite magnetic, clearly containing significant magnetite. While magnetite spinel is generally absent in the critical zone, including the chromitites, magnetite is common in the Platreef \citep{kinnaird2005}, the basement-contaminated PGE-rich north limb corollary of the critical zone’s Merensky Reef. Because the upper critical zone at Kliprivier Farm also formed at the chamber floor, it may be similar in character and PGE grade to the Platreef \citep{vandermerwe2007}.}{66}{Chapter5/Figures/generated/field_photos/DrBV66.jpg}{\textbf{Photo 66.} Massive ($\sim$20 cm) non-magnetic chromitite with deformed blocks of anorthosite, preferentially lying along the bedding planes of the area. The lens cap is 5.5 cm. Anorthosite blocks appear to have been ductilely deformed, perhaps by entrainment within a hot chromite slurry. Like the chromitites hosted within the anorthosites, there are oikocrystic pyroxenes which appear to have crystallized late. The groundmass is chromitite with an abundant feldspar matrix, perhaps only 50\% cumulus chromite, 30\% interstitial plagioclase, 20\% oikocrystic pyroxenes. Particularly in the top right corner, one can see a striking increase in modal chromitite approaching the boundary of the entrained anorthosite block. Despite the last photo resembling UG-1 stringers, it is structurally above this massive chromitite (UG-1 stringers are below: \citealp{voordouw2009,mukherjee2017}). The orange color on the surface is lichen.}
\ChFiveWidePair{72}{Chapter5/Figures/generated/field_photos/DrBV72.jpg}{\textbf{Photo 72.} Blobby aggregates of pyroxene of the coarsely crystalline Disturbed Gabbronorite (Dcg) of the Kliprivier Monocline. The orange magnet is the size of a typical pen, $\sim$13 cm long. Could these be precursors to spheroidal aggregates (SAs)? Photo~\ChFivePhoto{72} shows a corona of pyroxene depletion surrounding the pyroxene aggregate.}{76}{Chapter5/Figures/generated/field_photos/DrBV76.jpg}{\textbf{Photo 76.} Decimetric xenolith of fine grained marginal zone rock within Disturbed Gabbronorite. The orange magnet is the size of a typical pen, $\sim$13 cm long. The xenolith exhibits strong resistance banding that is subtle or absent in the coarse grained disturbed gabbronoritic host, and the contacts between the two are sharp. Interestingly, aggregates of pyroxenes have glommed to the rims of the xenolith.}
\ChFiveWidePair{77}{Chapter5/Figures/generated/field_photos/DrBV77.jpg}{\textbf{Photo 77.} Decimetric xenolith of fine grained marginal zone rock within Disturbed Gabbronorite. The orange magnet is the size of a typical pen, $\sim$13 cm long. This xenolith has a comparatively diffuse contact with the disturbed gabbronoritic host relative to that of photo~\ChFivePhoto{76}, perhaps suggesting that it had been incorporated within the magma for longer, or had been reheated more after entrainment.}{86}{Chapter5/Figures/generated/field_photos/DrBV86.jpg}{\textbf{Photo 86.} A thin ($<$15 cm) band of spotted anorthosite in the mid-lower Main Zone, exposed in weathered and clean faces. The hammer head is about $\sim$5 cm in width. The spotted anorthosites were much harder to spot in the field than the 3 dramatic packages of mottled anorthosites found in the valley and further downsection, both above and below the Draaikraal Lineament. \citet{latypovetal2022} correlated this layer to the Upper Mottled Anorthosite (UMA), and assumed it was laterally continuous with the rocks of photos~\ChFivePhoto{49}, \ChFivePhoto{53}, and \ChFivePhoto{54} south of the Draaikraal Lineament. While I agree that photos~\ChFivePhoto{49}, \ChFivePhoto{53}, and \ChFivePhoto{54} show the UMA, I instead correlate the UMA above the lineament much further to the east (photo~\ChFivePhoto{90}), where I identified the characteristic fine-grained, bright green gabbronorite with blebby, anhedral phenocrystic ($<$2 cm) CPX that I had seen nowhere else but as footwall to the UMA. I also interpret photos~\ChFivePhoto{114} and \ChFivePhoto{117} as belonging to the UMA.}
\ChFiveWidePair{87}{Chapter5/Figures/generated/field_photos/DrBV87.jpg}{\textbf{Photo 87.} Similarly to photo~\ChFivePhoto{86}, this rock was identified as a “mottled anorthosite” by \citet{latypovetal2022} and was correlated to the MMA marker bed. The pen used for scale is about 1 cm in diameter. I instead refer to this unit again as a “spotted anorthosite” since it bears none of the striking features of the MMA seen in photos~\ChFivePhoto{93} and \ChFivePhoto{103}--\ChFivePhoto{109} exposed to the south of the Draaikraal Lineament. I argue that the true corollary to the MMA north of the Draaikraal Lineament is much further east, (see the thick mottled anorthosite package in photo~\ChFivePhoto{115}).}{90}{Chapter5/Figures/generated/field_photos/DrBV90.jpg}{\textbf{Photo 90.} Coarse mottled anorthosite, with a $\sim$10 cm wallet for scale. Unfortunately, I did not have the foresight to photograph the unique footwall to these anorthosites in place: the fine-grained, bright green gabbronorite with blebby, anhedral phenocrystic ($<$2 cm) CPX that I found only here and at Draaikraal (at the outcrop in photos~\ChFivePhoto{49}, \ChFivePhoto{53}, and \ChFivePhoto{54}) and in the float traced along strike north of that outcrop. I did sample it, however (D3SG).}
\ChFiveWidePair{95}{Chapter5/Figures/generated/field_photos/DrBV95.jpg}{\textbf{Photo 95.} This a particularly striking example of resistance banding in fine-grained marginal zone sill phase float. The hammer head is $\sim$10 cm across. This feature was immediately noticeable in the fine-grained magnetic marginal zone rocks, but I didn’t photograph it until later, assuming at first that the banding was a secondary, unimportant textural feature. The ubiquitousness of these bands in fine-grained marginal zone rocks and in xenoliths within the plutonic RLS. These bands are most often subparallel and vertical planes. Similar resistance banding is argued to form in sedimentary rocks by localized cataclasis and recrystallization along fault planes during brittle deformation, decreasing the porosity of the affected plane to become more resistant to later chemical weathering \citep{hill1990}. While it is tempting to ascribe the formation of these bands to subsolidus (or near-subsolidus) deformation during isostatic flexure of the cooling RLS, these bands’ formation must predate the intrusion of the Main Zone of the RLS, which includes xenoliths of fine-grained Marginal Zone rocks containing these bands (photo~\ChFivePhoto{91}). An SEM/petrographic investigation of these bands would be helpful to understand their origins and implications.}{96}{Chapter5/Figures/generated/field_photos/DrBV96.jpg}{\textbf{Photo 96.} Two $\sim$2 cm spherical pyroxenite aggregates found in a leucocratic Marginal Sill Phase outcrop, and a single $\sim$5 cm pyroxenite aggregate with a long ($\sim$10 cm), thin pyroxenite tail. The hammer’s handle is $\sim$4 cm wide, for scale. This feature is interpreted similarly to a diapir, implying that the pyroxenite aggregate was ductily deformable and migrating through a comparatively rigid/viscous crystal mush. The requirement that the pyroxenite aggregate be ductily deformable is inconsistent with the model of \citet{khoza2021}, in which preexisting and relatively cool pyroxenite layers are brecciated by an intruding (gabbro)noritic magma, but is consistent with the models of \citet{lee1979}, in which pyroxenite is preferentially wetted by a secondary, immiscible bridging liquid; and \citet{mitchell2023}, who argue that pyroxenite SAs are formed similarly to mafic enclaves in granitoids, by disaggregation of an incoming pyroxenitic magma within a pre-existing, comparatively viscous/rigid (gabbro)noritic host.}
\ChFiveWidePair{98}{Chapter5/Figures/generated/field_photos/DrBV98.jpg}{\textbf{Photo 98.} Anastomosing, non-planar fine-grained anorthositic (?) veins in the fine-grained Marginal Zone. Similar to veins observed in photo~\ChFivePhoto{6}. The hammer head is $\sim$10 cm across. Anorthositic veins include blocks of host fine-grained Marginal Zone rocks with indistinct, irregular contacts. Such veins may represent incompletely formed anorthositic layers formed by reactive porous flow of a magma or hydrothermal fluid sourced from dehydrating/melting country rocks. Note the presence of modal layering within the anorthosite (?) of photo~\ChFivePhoto{98}.}{100}{Chapter5/Figures/generated/field_photos/DrBV100.jpg}{\textbf{Photo 100.} Highly migmatized Houtenbek shales in weathered and fresh faces. The hammer head is $\sim$10 cm across. Deformed, fractured, and melted metasedimentary rocks swim in abundant, pink, leucosomal crystalline granite. Partial melt of crustal rocks was generated in-situ within the sedimentary floor rocks of the Bushveld, and it is able to effectively separate from the source rocks. Some amount of melts like these likely migrated upward to contaminate the plutonic RLS and Marginal Zone rocks. This contamination effect would be strongest for magmas near the intrusion floor.}
\ChFiveWidePair{109}{Chapter5/Figures/generated/field_photos/DrBV109.jpg}{\textbf{Photo 109.} Mottled anorthosite (MMA) south of the Draaikraal Lineament. The hammer head is $\sim$10 cm across. The interior of a pyroxenite slab separating two layers of mottled anorthosite (as pictured in photo~\ChFivePhoto{108}) is well exposed by high-angle fracture. The interiors of these pyroxenite slabs are finer grained and more anorthosite rich than the exteriors. In general, I observe coarse grains where mineral phases are well-separated, such as in mottled anorthosites or at anorthosite-pyroxenite interfaces.}{111}{Chapter5/Figures/generated/field_photos/DrBV111.jpg}{\textbf{Photo 111.} $\sim$ 12 cm bed of Spotted anorthosite north of the Draaikraal Lineament, upsection of the UMA. The hammer head is $\sim$10 cm across. A blocky, melanocratic xenolith/autolith is included near the middle of the bed, and there is another blobby pyroxene-rich blob mantled by pure anorthosite at the bottom of the bed. The spotty Pyroxenes do not appear oikocrystic. They are larger near the middle of the spotted anorthosite and generally absent near the bottom. The contacts with the host gabbronorite are indistinct and irregular.}
\ChFiveTallPair{29}{Chapter5/Figures/generated/field_photos/DrBV29.jpg}{\textbf{Photo 29.} A Well-foliated hornfelsized inclusion of Houtenbek shale/dolomite within the Marginal Sill Phase on the Kliprivier Synform. The foliation shows clear plastic deformation, indicative of synmagmatic (hot) deformation. While previous authors have argued that the Kliprivier synform preceded Bushveld magmatism \citep{vandermerwe2007}, this work best supports the interpretation that deformation the RLS margin on Kliprivier Farm was synmagmatic.}{53}{Chapter5/Figures/generated/field_photos/DrBV53.jpg}{\textbf{Photo 53.} The same UMA outcrop as pictured in 49. The UMA package includes a striking mottled anorthosite above a coarse feldspathic pyroxenite with cumulus textures and decimeter-scale anorthosite inclusions with indistinct contacts with the host pyroxenite. Most of the anorthositic inclusions within the pyroxenite are roughly situated in a horizon $\sim$15 cm above the base of the feldspathic pyroxenite layer, but one (especially indistinct) inclusion is near the upper contact and is nearly continuous with overlying mottled anorthosite. The two holes in the outcrop were likely from explosives for clearing the outcrop from the road, but may also have been from sampling from previous paleomagnetic or geochemical studies: they are $\sim$ 4 cm across for scale.}
